\chapter{Grundlagen}
\label{ch:02_grundlagen}

% Kapitelübersicht.
In diesem Kapitel werden die theoretischen Grundlagen der untersuchten Rendering-Pipelines und der verwendeten Verfahren
definiert und erörtert.
Hierzu zählen die technischen Grundlagen der diversen Rendering-Paradigmen, die Spezifikation der physischen
Hardware-Flaschenhälse sowie die Einordnung der Konzepte im Kontext der Stilisierung.


\section{Rendering-Paradigmen}
\label{sec:0201_rendering_paradigms}

% Einleitung.
Im Folgenden werden die formalen Definitionen und die Gegenüberstellung von klassischen und entkoppelten Render-Pfaden
unter Berücksichtigung ihrer inhärenten Skalierbarkeit erörtert.

\subsection{Forward Rendering}
\label{subsec:020101_forward_rendering}

% Funktionsweise.
Naives Forward Rendering (im Folgenden als Bezeichnung von Single-Lighting-Pass Forward-Varianten verwendet) zeichnet
sich durch das Iterieren über einzelne Geometrien und Bearbeiten von Shading-Operationen pro Objekt aus.
Es erfolgt zunächst eine Evaluierung aller Lichter, die ein potenzielles Objekt in einer Szene beeinflussen.
In einer weiteren Schleife werden diese Lichter anschließend für das jeweilige Objekt zusammen mit etwaigen anderen
Materialeigenschaften verrechnet und
angewendet~\cite[Folie~\textit{Forward Rendering – Single Pass}]{valientDeferredRenderingKillzone}
(vgl.\ ~\ref{fig:naive_forward_rendering_flow}).
Die Beleuchtungsberechnung finden also untrennbar gekoppelt im selben Render-Pass (Fragment-Shader) statt.

\begin{figure}[htbp]
    \centering
    \resizebox{0.85\textwidth}{!}{%
        \begin{tikzpicture}[
            >=Latex,
            node distance=0.7cm and 1.5cm,
            font=\sffamily\small,
            base/.style={draw, align=center, minimum height=0.7cm, minimum width=3.5cm},
            startstop/.style={base, rounded corners, fill=gray!15},
            process/.style={base, fill=white},
            decision/.style={draw, diamond, aspect=2, inner sep=2pt, text width=2.2cm, align=center}
        ]

            % Nodes
            \node (start) [startstop] {Szenen-Rendering starten};

            % Outer Loop: Objects
            \node (dec_obj) [decision, below=of start] {Alle Objekte\\gerendert?};
            \node (end) [startstop, right=of dec_obj, xshift=1cm] {Rendering abgeschlossen};
            \node (bind_obj) [process, below=of dec_obj] {Geometrie \& individuelles\\Material/Shader binden};

            % Lights
            \node (dec_light) [decision, below=of bind_obj] {Alle Lichter\\evaluiert?};
            \node (calc_light) [process, right=of dec_light, xshift=0.8cm] {Beleuchtung akkumulieren
            \\(Fragment-Shader)};

            % Draw Call
            \node (draw) [process, below=of dec_light] {Pixel in Framebuffer schreiben};

            % Edges - Main Flow
            \draw[->] (start) -- (dec_obj);
            \draw[->] (dec_obj) -- node[above, font=\footnotesize] {Ja} (end);
            \draw[->] (dec_obj) -- node[right, font=\footnotesize] {Nein} (bind_obj);
            \draw[->] (bind_obj) -- (dec_light);
            \draw[->] (dec_light) -- node[right, font=\footnotesize] {Ja} (draw);

            % Edges - Lights
            \draw[->] (dec_light) -- node[above, font=\footnotesize] {Nein} (calc_light);
            \draw[->] (calc_light.north) -- ++(0,0.5) -| (dec_light.north);

            % Edges - Objects
            \draw[->] (draw.west) -- ++(-1.0,0) |- node[near start, left, font=\footnotesize] {Nächstes Objekt}
            (dec_obj.west);

        \end{tikzpicture}
    }
    \caption{UML-Aktivitätsdiagramm des naiven Forward-Rendering-Paradigmas.
    Die direkte Verschachtelung der Objekt- und Lichtevaluierung zeigt die inhärente quadratische Komplexität.}
    \label{fig:naive_forward_rendering_flow}
\end{figure}

% Komplexität.
Dieser Umstand führt zu einem grundlegenden theoretischen Aufwand des Forward-Renderings, der sich proportional zur
Anzahl der Geometrien und der einfallenden Lichtquellen verhält (quadratische Komplexität).
~\cite[Abschn.~1]{Thaler2011DeferredRendering}.
\begin{equation}
    \mathcal{O}(N_{\text{obj}} \cdot N_{\text{light}})
    \label{eq:forward_shading_complexity}
\end{equation}
Wobei $N_{\text{obj}}$ die Objektanzahl und $N_{\text{light}}$ die Anzahl der aktiven, Lichtquellen definiert.

% Overdraw-Teaser.
Auf Hardware-Ebene manifestiert sich diese Komplexität in der Rasterisierung.
Da die Beleuchtungsschleife für jedes erzeugte Fragment ausgeführt wird, skaliert die Rechenlast massiv mit der
Pixelfläche der Objekte~\cite[Abschn.~3.2]{haradaForwardBringingDeferred2012}.
Verdeckt ein nachfolgend gezeichnetes Objekt ein bereits berechnetes Fragment, war die vorherige, rechenintensive
Lichtauswertung obsolet.
Dieses ineffiziente Verhalten wird als \textit{Shading Overdraw} bezeichnet und in
Abschnitt~\ref{2.2.1} detailliert beleuchtet.

% Vorteile.
Die direkte Kopplung von Geometrie und Material bietet in der Praxis jedoch auch diverse Vorteile.
Es kann jedes Objekt ohne sonstige Workarounds seinen eigenen hochspezialisierten Shader
binden~\cite[Abschn.~1]{haradaForwardBringingDeferred2012}.
Dies ist insbesondere vor dem Hintergrund der Evaluation dieser Arbeit bezüglich selektiver Stilisierungen mittels
\ac{NPR}-Effekten von Relevanz.

\subsection{Deferred Rendering}
\label{subsec:020102_deferred_rendering}

% Funktionsweise.
Deferred-Rendering zeichnet sich im Gegensatz zu Forward-Rendering durch die eindeutige Trennung von
Geometrieverarbeitung und dem Shading der Geometrie aus.
Die Geometrien sowie alle damit assoziierten Werte, wie etwa Normalen (Ausrichtungen), Albedos
(Farbinformationen) oder Tiefeninformationen, werden in einem eigenständigen, separaten Pass verarbeitet und in
einen Zwischenspeicher, dem sogenannten \ac{G-Buffer}, in Form eines \ac{FBO} geschrieben (bildschirmfüllendes Quad).
In diesem Pass werden explizit keine Shading-Operationen, wie etwa die Beleuchtung der Objekte, durchgeführt.
Erst in einem nachgelagerten Pass werden schließlich dei IInformation aus dem \ac{G-Buffer} verwendet, um die
Beleuchtung und sonstige Materialeigenschaften auf Screen-Space-Ebene auszuwerten und zu berechnen
(vgl.\ ~\ref{fig:deferred_rendering_flow}).~\cites[Abschn.~2]{Thaler2011DeferredRendering}
[Kap.~20.1]{akenine-mollerRealtimeRendering2019}
[Folie~\textit{Deferred Rendering}]{valientDeferredRenderingKillzone}

% Diagram.
\begin{figure}[htbp]
    \centering
    \resizebox{0.95\textwidth}{!}{%
        \begin{tikzpicture}[
            >=Latex,
            node distance=0.8cm and 1.5cm,
            font=\sffamily\small,
            base/.style={draw, align=center, minimum height=0.7cm, minimum width=3.5cm},
            startstop/.style={base, rounded corners, fill=gray!15},
            process/.style={base, fill=white},
            decision/.style={draw, diamond, aspect=2, inner sep=2pt, text width=2.2cm, align=center},
            storage/.style={base, fill=blue!05, rounded corners=2mm}
        ]

            % Nodes
            \node (start) [startstop] {Szenen-Rendering starten};

            % Geometry Pass
            \node (dec_obj) [decision, below=of start] {Alle Objekte\\verarbeitet?};
            \node (write_g) [process, below=of dec_obj] {Geometrie, Normalen, Albedo\\extrahieren};
            \node (store_g) [process, below=of write_g] {In G-Buffer (FBO) schreiben\\(Explizit kein Shading)};

            % Buffer State Transition
            \node (gbuffer) [storage, right=of dec_obj, xshift=1.5cm] {G-Buffer vollständig\\(Geometrie-Pass beendet)};

            % Lighting Pass
            \node (dec_light) [decision, below=of gbuffer] {Alle Lichter\\evaluiert?};
            \node (read_g) [process, below=of dec_light] {Daten aus G-Buffer lesen\\(Screen-Space)};
            \node (calc_light) [process, below=of read_g] {Shading berechnen\\(Nur sichtbare Fragmente)};

            % End
            \node (end) [startstop, right=of dec_light, xshift=1.5cm] {Rendering abgeschlossen};

            % Edges - Geometry Loop
            \draw[->] (start) -- (dec_obj);
            \draw[->] (dec_obj) -- node[right, font=\footnotesize] {Nein} (write_g);
            \draw[->] (write_g) -- (store_g);
            \draw[->] (store_g.west) -- ++(-1.0,0) |- node[near start, left, font=\footnotesize] {Nächstes Objekt}
            (dec_obj.west);

            % Edges - Transition
            \draw[->] (dec_obj) -- node[above, font=\footnotesize] {Ja} (gbuffer);
            \draw[->] (gbuffer) -- (dec_light);

            % Edges - Lighting Loop
            \draw[->] (dec_light) -- node[right, font=\footnotesize] {Nein} (read_g);
            \draw[->] (read_g) -- (calc_light);
            \draw[->] (calc_light.west) -- ++(-1.0,0) node[below right, font=\footnotesize, xshift=-2cm]
                {Nächstes Licht} |- (dec_light.west);

            % Edges - End
            \draw[->] (dec_light) -- node[above, font=\footnotesize] {Ja} (end);

        \end{tikzpicture}
    }
    \caption{UML-Aktivitätsdiagramm des Deferred-Rendering-Paradigmas.
    Die algorithmische Entkopplung von Geometrieverarbeitung (linker Pfad) und Beleuchtungsberechnung
        (mittlerer Pfad) zeigt die inhärent lineare Komplexität.
        Shading-Operationen erfolgen ausschließlich im Screen-Space für sichtbare Fragmente.}
    \label{fig:deferred_rendering_flow}
\end{figure}

% Komplexität.
Algorithmisch bedingt die additive Entkopplung demnach einen linearen, theoretischen Aufwand für
Deferred-Rendering-Paradigmen~\cite[Abschn.~3]{fonsecaDeferredShadingTutorial}.
\begin{equation}
    \mathcal{O}(N_{\text{obj}} + N_{\text{light}})
    \label{eq:deferred_shading_complexity}
\end{equation}
Hierbei definiert $N_{\text{obj}}$ die Anzahl der Objekte während der initialen Geometriephase (G-Buffer-Pass) und
$N_{\text{light}}$ die Anzahl der aktiven Lichtquellen.
Die Shading-Operationen werden folglich nur für die tatsächlich sichtbaren Fragmente (Pixel) berechnet.

% Uber Shader/
Werden Deferred-Ansätze nur in einem einzigen, dem Geometrie-Pass nachgelagerten, Pass umgesetzt, kommt es bei
der Umsetzung unterschiedlicher Materialeigenschaften und möglichen Stilisierungen zwangsweise zu einer monolithischen
Shader-Struktur im finalen Fragment-Shader und sogenannten
\textit{Uber-Shadern}~\cites[S.~881]{akenine-mollerRealtimeRendering2019}[Abschn.~2.6]{Thaler2011DeferredRendering}
[Abschn.~6]{schiedDeferredAttributeInterpolation2015}.
Für die in dieser Arbeit implementierten Single-Pass Deferred-Pipeline
(im weiteren Verlauf \textit{Deferred-Uber} genannt) ist dieser Umstand vor dem Hintergrund der selektiven stilisierung
mit \ac{NPR}-Effekten besonders bedeutend.

% Light Volumes.
Ein erweiterter und gängiger Ansatz der Implementierung von Deferred-Pipelines stellt eine Multi-Pass Shading
Architektur unter der verwendung von \textit{Light-Volumes} dar~\cites[Abschn.~2]{mallettDeferredAdaptiveCompute2018}
[Folie~"Conventional Deferred Shading"]{lauritzenDeferredRenderingCurrent}.
In diesem Ansatz werden die Lichtquellen um spezielle Proxy-Geometrien erweitert, die jeweils den Wirkungsbereich dieser
Lichtquellen darstellen~\cite[Abschn.~2.2]{Thaler2011DeferredRendering}.
Für diesen Ansatz muss das eigentliche Shading in zwei separate Pässe aufgeteilt werden.
Nach dem initialen Geometrie-Pass, in dem der \ac{G-Buffer} ohne die entsprechenden Proxy-Geometrien beschrieben wurde,
werden die Light-Volumes in einem separaten Lighting-Pass rasterisiert~\cites[Abschn.~2]{Thaler2011DeferredRendering}
[S.~884--886]{akenine-mollerRealtimeRendering2019}.
Innerhalb dieses Prozesses werden ausschließlich die Beleuchtungseinflüsse für die betroffenen Fragmente durch das
Auslesen des \ac{G-Buffer}s berechnet und in ein separates \ac{FBO} (\textquote{Lighting-Buffer})
geschrieben~\cite[Abschn.~3.1]{fonsecaDeferredShadingTutorial}.
Die Einflüsse etwaiger überlappender Light-Volumes werden durch automatisierte additive \textit{Blending}-Operationen
auf den speziell dafür ausgelegten \acp{ROP} der \ac{GPU}
akkumuliert~\cite[Kap.~11.2.4.10]{gregoryGameEngineArchitecture2018}.
In einem finalen Resolve-Pass werden schließlich die Basis-Materialeigenschaften (wie etwa Farbinformationen) aus dem
\ac{G-Buffer} mit der akkumulierten Beleuchtung kombiniert und um finale Shading- oder Post-Processing-Effekte
ergänzt~\cite[Abschn.~2.3]{Thaler2011DeferredRendering}.

\begin{figure}[htbp]
    \centering
    \resizebox{0.95\textwidth}{!}{%
        \begin{tikzpicture}[
            >=Latex,
            node distance=0.8cm and 1.5cm,
            font=\sffamily\small,
            base/.style={draw, align=center, minimum height=0.7cm, minimum width=3.5cm},
            startstop/.style={base, rounded corners, fill=gray!15},
            process/.style={base, fill=white},
            decision/.style={draw, diamond, aspect=2, inner sep=2pt, text width=2.2cm, align=center},
            storage/.style={base, fill=blue!05, rounded corners=2mm}
        ]

            % Nodes
            \node (start) [storage] {G-Buffer initialisiert\\(Geometrie-Pass beendet)};

            % Lighting Pass
            \node (dec_vol) [decision, below=of start] {Alle Light-Volumes\\gerastert?};
            \node (raster) [process, below=of dec_vol] {Proxy-Geometrie\\rasterisieren};
            \node (read_g) [process, below=of raster] {G-Buffer auslesen\\(Nur betroffene Fragmente)};
            \node (accumulate) [process, below=of read_g] {In Lighting-Buffer schreiben\\(Additives Blending
            via \ac{ROP}s)};

            % Resolve Pass
            \node (buffer_ready) [storage, right=of dec_vol, xshift=1.8cm] {Lighting-Buffer\\vollständig};
            \node (resolve) [process, below=of buffer_ready] {Resolve-Pass:\\Albedo \& Licht kombinieren};
            \node (post) [process, below=of resolve] {Finales Shading \&\\Post-Processing};
            \node (end) [startstop, below=of post] {Rendering abgeschlossen};

            % Edges - Lighting Pass
            \draw[->] (start) -- (dec_vol);
            \draw[->] (dec_vol) -- node[right, font=\footnotesize] {Nein} (raster);
            \draw[->] (raster) -- (read_g);
            \draw[->] (read_g) -- (accumulate);
            \draw[->] (accumulate.west) -- ++(-1.0,0) |- node[pos=0.5, above left, font=\footnotesize]
                {Nächstes Volume} (dec_vol.west);

            % Edges - Resolve Pass
            \draw[->] (dec_vol) -- node[above, font=\footnotesize] {Ja} (buffer_ready);
            \draw[->] (buffer_ready) -- (resolve);
            \draw[->] (resolve) -- (post);
            \draw[->] (post) -- (end);

        \end{tikzpicture}
    }
    \caption{UML-Aktivitätsdiagramm der Multi-Pass-Architektur für Deferred-Rendering mittels Light-Volumes.
    Im Lighting-Pass (linker Pfad) werden Proxy-Geometrien gerastert und überlappende Einflüsse per Hardware-Blending
    in einem separatem FBO akkumuliert.
    Der Resolve-Pass (rechter Pfad) kombiniert diese Daten anschließend mit den initialen Materialinformationen.}
    \label{fig:deferred_light_volumes_flow}
\end{figure}


%\subsection{Deferred Rendering}
%\label{subsec:020102_deferred_rendering_core}
%Erklärung der entkoppelten Pipeline und der Zwischenspeicherung von Geometriedaten in separaten Puffern. Aufzeigen der Effizienzgewinne bei der Schattierung durch den Einsatz geometrischer Hilfskörper.
%
%
%\section{Hardware-Limitierungen und Performance-Metriken}
%\label{sec:0202_hardware_limitations_and_metrics}
%
%Dieser Abschnitt spezifiziert die physischen Ausführungseinheiten moderner Grafikhardware, welche die kritischen Engpässe bei der praktischen Ausführung der Render-Pfade bilden.
%
%\subsection{Fragment-Overdraw und Fillrate}
%\label{subsec:020201_fragment_overdraw_and_fillrate}
%Thematisierung mehrfach gezeichneter Pixel auf identischen Koordinaten. Betrachtung der resultierenden Last für die Füllrate und softwareseitiger Optimierungen wie dem Z-Prepass.
%
%\subsection{Speicherbandbreite (Memory Bandwidth)}
%\label{subsec:020202_memory_bandwidth_limitations}
%Analyse der Durchsatzgrenzen beim Zugriff auf den Videospeicher. Erläuterung des Lese- und Schreibaufwands durch breite Pufferstrukturen und additive Mischvorgänge.
%
%\subsection{Register Pressure}
%\label{subsec:020203_register_pressure_limitations}
%Behandlung der limitierten internen Speichereinheiten pro Rechenkern. Beschreibung der Auswirkungen auf die Thread-Parallelisierung bei der Nutzung großer, kombinierter Shader-Programme.
%
%\subsection{Dynamic Branching und Warp Divergence}
%\label{subsec:020204_dynamic_branching_and_warp_divergence}
%Erläuterung der synchronen Funktionsweise von Thread-Gruppen auf Hardware-Ebene. Beschreibung der sequenziellen Serialisierung bei ungleichen Verzweigungsbedingungen auf Pixelebene.
%
%
%\section{Grundlagen der NPR Pipeline}
%\label{sec:0203_non_photorealistic_rendering_fundamentals}
%
%Dieser Abschnitt kategorisiert die raumbezogenen Konzepte der NPR-Generierung zur selektiven, isolierten Stilisierung von Szenenobjekten.
%
%Abwägung von Stilisierungsverfahren direkt auf der dreidimensionalen Szenengeometrie im Vergleich zu nachgelagerten, globalen Filtern auf dem fertigen zweidimensionalen Bild. Gegenüberstellung von präziser Objekttrennung und Laufzeit-Konstanz.
